La solubilité d'un soluté dans un solvant est la masse maximale de ce soluté que
l'on peut dissoudre dans un litre de solution. Elle s'exprime en
$\unit{\g\per\L}$. Elle dépend de la température. Lorsque la solubilité est
atteinte la solution est qualifiée de saturée.

Le tableau ci-dessous représente la solubilité dans l'eau de quelques composés
ioniques à $\qty{20}{\degreeCelsius}$.

\begin{minipage}{\linewidth}
    \begin{table}[H]
    \centering
    \setlength{\arrayrulewidth}{0.8pt}
    \renewcommand{\arraystretch}{1.5}
    \begin{tabularx}{\textwidth}{|C|C|}
        \hline
        Composés & Solubilité à $\qty{20}{\degreeCelsius}$~(\unit{\g\per\L}) \\
        \hline
        \ce{NaCl} & 360 \\
        \hline
        \ce{NH4Cl} & 372 \\
        \hline
        \ce{KNO3} & 316 \\
        \hline
    \end{tabularx}
\end{table}
\end{minipage}

\begin{enumerate}
    \item Écrire les équations des réactions de dissolution des solutés dans
          l'eau.
    \item Calculer la concentration molaire des ions présents en solution.
\end{enumerate}

